% v10, May 27 2026, Claude Opus 4.7

\documentclass[11pt]{article}

\usepackage[margin=1in]{geometry}
\usepackage{amsmath, amssymb, amsthm}
\usepackage{mathtools}
\usepackage{xcolor}
\usepackage{hyperref}
\hypersetup{colorlinks=true,
            linkcolor=blue!50!black,
            citecolor=blue!50!black,
            urlcolor=blue!50!black}

%---- math macros ----
\newcommand{\R}{\mathbb{R}}
\newcommand{\E}{\mathbb{E}}
\renewcommand{\Pr}{\mathbb{P}}
\newcommand{\eps}{\varepsilon}
\newcommand{\Tr}{\operatorname{Tr}}
\newcommand{\Trs}{\operatorname{Tr}_{\sigma}}
\newcommand{\diag}{\operatorname{diag}}
\newcommand{\opnorm}[1]{\left\lVert #1 \right\rVert}
\newcommand{\del}{\delta}
\DeclareMathOperator{\Span}{span}
\DeclareMathOperator{\im}{im}
\DeclareMathOperator{\rank}{rank}

\newcommand{\tL}{\widetilde{L}}
\newcommand{\Lt}{L^{\dagger}}
\newcommand{\Lth}{L^{\dagger/2}}

%---- theorem environments ----
\newtheorem{theorem}{Theorem}
\newtheorem{lemma}[theorem]{Lemma}
\newtheorem{proposition}[theorem]{Proposition}
\newtheorem{claim}[theorem]{Claim}

\title{$\eps$-Light Subsets in Graphs}
\author{}
\date{May 27, 2026}

\begin{document}
\maketitle

\section*{Statement and Answer}

Let $G=(V,E,w)$ be a weighted graph on $n=|V|$ vertices with Laplacian
$L=\sum_{(s,t)\in E}w(s,t)\,(\delta_s-\delta_t)(\delta_s-\delta_t)^{\!\top}$.
For $S\subseteq V$ let $G_S=(V,E(S,S))$ keep only edges with both endpoints in
$S$ and let $L_S$ be its Laplacian (same ambient dimension as $L$). $S$ is
\emph{$\eps$-light} iff $\eps L - L_S\succeq 0$.

\medskip
\noindent\textbf{Theorem.} \emph{For every weighted graph $G=(V,E,w)$ and every
$\eps\in(0,1)$, there exists an $\eps$-light subset $S\subseteq V$ with
$|S|\ge \lfloor \eps n/42\rfloor$.}

\medskip
\noindent The answer is \textbf{YES}, with the explicit constant $c=1/42$.

\section*{Notation and basic facts}

Write $\Lt$ for the Moore--Penrose pseudoinverse and $\Lth$ for its PSD square
root. Define
\[
  \tL_S      := \Lth L_S \Lth, \qquad
  \tL_{S,T}  := \Lth L_{S,T}  \Lth, \qquad
  \tL_{S,t}  := \Lth L_{S,\{t\}}\Lth.
\]
Then $\Lth L \Lth = \Pi$, the orthogonal projector onto $\Span\{L\}$, and the
condition $\eps L\succeq L_S$ is equivalent (by conjugation with $\Lth$ on
$\im L$) to the operator-norm bound
\begin{equation}\label{eq:opnorm-form}
  \opnorm{\tL_S}_{\mathrm{op}} \;\le\; \eps.
\end{equation}

\medskip
\noindent\emph{Leverage scores.}\quad The leverage score of an edge $(s,t)$ is
\[
  \ell(s,t)
   := w(s,t)\,(\delta_s-\delta_t)^{\!\top}\Lt(\delta_s-\delta_t)
    = \Tr\!\bigl(L_{\{s\},\{t\}}\Lt\bigr),
\]
and zero on non-edges. For $S,T\subseteq V$ set
$\ell(S,T):=\sum_{s\in S,t\in T}\ell(s,t)$ and write
$\ell(S):=\ell(S,V\setminus S)$ for the boundary leverage. By linearity
$\ell(S,T)=\Tr(\tL_{S,T})$. Foster's identity, in the form
$\sum_{(s,t)\in E}\ell(s,t)=n-1$ for connected $G$, follows from
$\sum_{(s,t)}L_{\{s\},\{t\}}=L$ and $\Tr(L\Lt)=\rank L = n-1$.

\medskip
\noindent\emph{Bounded-rank upper barrier.}\quad For a symmetric $A$ with
eigenvalues $\lambda_1\ge\cdots\ge\lambda_n$ and a real $u$, define
\[
  \Phi^u_{\sigma}(A) \;:=\;
  \begin{cases}
    \displaystyle\sum_{i=1}^{\sigma}\dfrac{1}{u-\lambda_i} & \text{if } u>\lambda_1,\\[2pt]
    \infty & \text{otherwise,}
  \end{cases}
\]
and the top-$\sigma$ trace $\Trs(A):=\sum_{i=1}^{\sigma}\lambda_i$. We write
$\Phi^u_\sigma(S)$ for $\Phi^u_\sigma(\tL_S)$, and note
$\Phi^u_\sigma(A)=\Trs\bigl((uI-A)^{-1}\bigr)$ when $u>\lambda_1$.

For the rest of this note we fix
\[
  \del := \frac{21}{n}, \qquad
  \phi := \frac{n}{21}, \qquad
  \sigma := \lfloor \eps n/42 \rfloor.
\]

\section*{Proof}

We build $S$ greedily, adding one vertex at a time and increasing the barrier
parameter $u$ by $\del$ at each step. Throughout, two invariants are
maintained: the bounded-rank barrier stays at most $\phi$, and the boundary
leverage stays at most $4|S|$.

\subsection*{Step 1: Initial set and target}

Pick any $v_0\in V$ and set $S_0:=\{v_0\}$, $u_0:=\eps/2$. Since $G_{S_0}$ has
no edges, $\tL_{S_0}=0$, and
\[
  \Phi^{u_0}_\sigma(S_0) \;=\; \frac{\sigma}{u_0}
   \;\le\; \frac{\eps n/42}{\eps/2} \;=\; \frac{n}{21} \;=\; \phi.
\]
After $\sigma$ greedy steps, the barrier will sit at $u_\sigma=u_0+\sigma\del$.
We chose parameters so that
\[
  u_\sigma = \tfrac{\eps}{2} + \sigma\cdot\tfrac{21}{n}
           \le \tfrac{\eps}{2} + \tfrac{\eps}{2} = \eps,
\]
so once we reach $|S|=\sigma+1$ with $\Phi^{u_\sigma}_\sigma(S)$ finite, every
eigenvalue of $\tL_S$ lies below $\eps$ and \eqref{eq:opnorm-form} gives
$\eps L\succeq L_S$.

\subsection*{Step 2: Greedy lemma}

\begin{lemma}\label{lem:greedy}
Suppose $|S|\le\sigma$, $\ell(S)\le 4|S|$, and $\Phi^u_\sigma(S)\le\phi$. Then
there exists $t\notin S$ with
\[
  \Phi^{u+\del}_\sigma(S\cup\{t\})\le\phi
  \qquad\text{and}\qquad
  \ell(S\cup\{t\})\le\ell(S)+4.
\]
\end{lemma}

\noindent The lemma is proved by two averaging arguments: more than half the
candidates $t$ keep the leverage budget (Step 3), and at least half keep the
barrier (Step 4). The two sets therefore intersect.

\subsection*{Step 3: Leverage-score averaging}

\begin{claim}\label{cl:Tlev}
For every $T\subseteq V$, $\sum_{t\in T}\ell(t,T)\le 2(|T|-1)$.
\end{claim}
\begin{proof}
By symmetry, $\sum_{t\in T}\ell(t,T)=2\Tr(L_T\Lt)$. Since $L_T\preceq L$, all
eigenvalues of $L_T\Lt$ lie in $[0,1]$, and $L_T$ has rank at most $|T|-1$, so
$\Tr(L_T\Lt)\le|T|-1$.
\end{proof}

For $t\notin S$,
\[
  \ell(S\cup\{t\}) \le \ell(S\cup\{t\},V\setminus S)
                   = \ell(S)+\ell(t,V\setminus S),
\]
so it suffices to control $\ell(t,V\setminus S)$. Claim~\ref{cl:Tlev} applied
to $T=V\setminus S$ gives
$\sum_{t\notin S}\ell(t,V\setminus S)<2|V\setminus S|$, whence by Markov more
than half of $t\notin S$ satisfy $\ell(t,V\setminus S)\le 4$.

\subsection*{Step 4: Barrier averaging}

This is the technical heart. The barrier-update lemma (Lemma~\ref{lem:bu}
below) reduces our task to showing that for at least half of $t\notin S$,
\begin{equation}\label{eq:U}
  U(S,t) \;:=\;
  \frac{\Trs(M^{-2}\tL_{S,t})}{\Phi^u_\sigma(S)-\Phi^{u+\del}_\sigma(S)}
   + \Trs(M^{-1}\tL_{S,t}) \;<\; 1,
\end{equation}
where $M:=(u+\del)I-\tL_S$. We claim
\begin{equation}\label{eq:total}
  \sum_{t\notin S} U(S,t) \;\le\; \tfrac{5}{\del} + 5\phi.
\end{equation}
Granting \eqref{eq:total}, since each $U(S,t)\ge 0$, more than half of
$t\notin S$ satisfy
\[
  U(S,t)\le\frac{2}{n-|S|}\!\left(\tfrac{5}{\del}+5\phi\right)
        \le \frac{2}{41n/42}\cdot\tfrac{10n}{21} < 1
\]
for $|S|\le\sigma\le n/42$.

\medskip
\noindent\emph{Proof of \eqref{eq:total}.}\quad Let $\Pi_S$ project onto
$\im\tL_S$, and $\Pi_S^\perp:=I-\Pi_S$. Since
$M=(u+\del)(\Pi_S+\Pi_S^\perp)-\tL_S$ and $\Pi_S^\perp\tL_S=0$,
\[
  M^p \;=\; (u+\del)^p\Pi_S^\perp + \bigl((u+\del)\Pi_S-\tL_S\bigr)^p,
  \qquad p\in\{-1,-2\}.
\]
By Ky Fan subadditivity of $\Trs$ \cite[Thm.~4.3.47a]{HJ12},
\begin{equation}\label{eq:split}
  \Trs(M^p\tL_{S,t}) \;\le\;
  \Trs\!\bigl((u+\del)^p\Pi_S^\perp\tL_{S,t}\bigr)
  + \Trs\!\Bigl(((u+\del)\Pi_S-\tL_S)^p\tL_{S,t}\Bigr).
\end{equation}

\smallskip
\emph{$\Pi_S^\perp$ piece.}\quad
$\Trs(\Pi_S^\perp\tL_{S,t})\le\Tr(\tL_{S,t})=\ell(S,t)$ (the PSD product has
non-negative eigenvalues), and summing,
\[
  \sum_{t\notin S}\Trs\!\bigl((u+\del)^p\Pi_S^\perp\tL_{S,t}\bigr)
   \le (u+\del)^p\,\ell(S)\le (u+\del)^p\cdot 4|S|.
\]

\smallskip
\emph{$\Pi_S$ piece.}\quad
$\sum_{t\notin S}\tL_{S,t}=\tL_{S,V\setminus S}\preceq\Pi$, and
$((u+\del)\Pi_S-\tL_S)^p$ is PSD on $\im\Pi_S$, so
\[
\begin{aligned}
\sum_{t\notin S}\Trs\!\Bigl(((u+\del)\Pi_S-\tL_S)^p\tL_{S,t}\Bigr)
  &\le \Tr\!\Bigl(((u+\del)\Pi_S-\tL_S)^p\tL_{S,V\setminus S}\Bigr)\\
  &\le \Tr\!\Bigl(((u+\del)\Pi_S-\tL_S)^p\Bigr)
   = \Tr(\Pi_S M^p\Pi_S) \le \Trs(M^p),
\end{aligned}
\]
the last step by the Ky Fan max principle.

\smallskip
Since all eigenvalues of $M$ are at most $u+\del$, for $p<0$ all eigenvalues of
$M^p$ are at least $(u+\del)^p$, so $\Trs(M^p)\ge\sigma(u+\del)^p\ge|S|(u+\del)^p$.
Combining,
\[
  \sum_{t\notin S}\Trs(M^p\tL_{S,t}) \;\le\; 4|S|(u+\del)^p + \Trs(M^p)
   \;\le\; 5\,\Trs(M^p).
\]
Now we use (Claim~\ref{cl:phi-diff})
$\Phi^u_\sigma(S)-\Phi^{u+\del}_\sigma(S)\ge\del\,\Trs(M^{-2})$ to bound the
first term of $U(S,t)$, and $\Trs(M^{-1})=\Phi^u_\sigma(S)\le\phi$ for the
second:
\[
  \sum_{t\notin S}U(S,t)
   \;\le\; \frac{5\,\Trs(M^{-2})}{\del\,\Trs(M^{-2})} + 5\,\Phi^u_\sigma(S)
   \;\le\; \tfrac{5}{\del}+5\phi,
\]
which is \eqref{eq:total}.\hfill$\square$

\subsection*{Step 5: Combining and conclusion}

Iterating Lemma~\ref{lem:greedy} $\sigma$ times starting from $S_0$ yields
$S$ with $|S|=\sigma+1$, $\ell(S)\le 4\sigma$, and
$\Phi^{u_\sigma}_\sigma(S)<\infty$, hence (by Step~1) every eigenvalue of
$\tL_S$ is at most $u_\sigma\le\eps$, so $\eps L\succeq L_S$. This proves the
theorem with $|S|\ge\lfloor\eps n/42\rfloor$.\qed

\section*{Auxiliary lemmas}

\begin{lemma}[Barrier update]\label{lem:bu}
Let $A,B\succeq 0$ and $\del>0$. Set $M:=(u+\del)I-A$. If
$\Phi^u_\sigma(A)<\infty$ and
\[
  \frac{\Trs(M^{-2}B)}{\Phi^u_\sigma(A)-\Phi^{u+\del}_\sigma(A)}
   + \Trs(M^{-1}B) \;<\; 1,
\]
then $\Phi^{u+\del}_\sigma(A+B)\le\Phi^u_\sigma(A)$.
\end{lemma}

\begin{proof}
$\Phi^u_\sigma(A)<\infty$ gives $u>\lambda_{\max}(A)$, hence
$M,M^{-1},M^{-2}\succ 0$. Write $B=CC^{\!\top}$. By Sherman--Morrison--Woodbury,
\[
  (M-CC^{\!\top})^{-1}
   = M^{-1} + M^{-1}C(I-C^{\!\top}M^{-1}C)^{-1}C^{\!\top}M^{-1}.
\]
Since $\Trs(M^{-1}B)=\Trs(C^{\!\top}M^{-1}C)<1$ the inverse exists and the
second summand is PSD, so $\lambda_{\max}(A+B)<u+\del$. By Ky Fan
subadditivity \cite[Thm.~4.3.47a]{HJ12},
\[
  \Phi^{u+\del}_\sigma(A+B)
   \le \Trs(M^{-1})
       + \Trs\!\bigl((I-C^{\!\top}M^{-1}C)^{-1}C^{\!\top}M^{-2}C\bigr).
\]
Since $\opnorm{C^{\!\top}M^{-1}C}\le\Trs(C^{\!\top}M^{-1}C)<1$, the operator
$(I-C^{\!\top}M^{-1}C)^{-1}\preceq(1-\Trs(C^{\!\top}M^{-1}C))^{-1}I$, and
monotonicity of $\Trs$ on PSD-ordered products gives
\[
  \Trs\!\bigl((I-C^{\!\top}M^{-1}C)^{-1}C^{\!\top}M^{-2}C\bigr)
   \le \frac{\Trs(M^{-2}B)}{1-\Trs(M^{-1}B)}.
\]
Combining with
$\Trs(M^{-1})=\Phi^u_\sigma(A)-(\Phi^u_\sigma(A)-\Phi^{u+\del}_\sigma(A))$
and the hypothesis gives the claimed bound.
\end{proof}

\begin{claim}\label{cl:phi-diff}
$\Phi^u_\sigma(S)-\Phi^{u+\del}_\sigma(S)\ge\del\,\Trs(M^{-2})$.
\end{claim}
\begin{proof}
For each top-$\sigma$ eigenvalue $\lambda_i$ of $\tL_S$,
\[
  \tfrac{1}{u-\lambda_i}-\tfrac{1}{u+\del-\lambda_i}
   = \tfrac{\del}{(u-\lambda_i)(u+\del-\lambda_i)}
   \ge \tfrac{\del}{(u+\del-\lambda_i)^2};
\]
sum over $i$.
\end{proof}

\section*{Remarks}

\begin{enumerate}\setlength\itemsep{4pt}
\item \textbf{Constant.} The constant $1/42$ comes from the parameter choices
$\del=21/n$, $\phi=n/21$, which balance the $5/\del$ and $5\phi$ contributions
in \eqref{eq:total}. It is not believed optimal; improving it is open.
\item \textbf{Key innovation.} The standard upper-barrier function tracks all
eigenvalues; here we track only the top $\sigma$. This adaptation, together
with the Ky Fan inequalities for $\Trs$, is what makes the argument go
through.
\item \textbf{Connection to spectral sparsification.} An $\eps$-light $S$ with
$|S|\ge c\eps n$ says the induced subgraph on $S$ contributes at most $\eps L$
spectrally, giving a perspective on ``light'' as ``spectrally independent''.
\item \textbf{Vertex vs.\ edge sampling.} The edge-partitioning analogue
admits an optimal constant of $1/2$ via the Kadison--Singer resolution. The
vertex-light variant is more delicate because edges sharing a vertex are
correlated; the deterministic barrier method bypasses concentration entirely.
\end{enumerate}

\begin{thebibliography}{9}
\bibitem{HJ12} R.~A.~Horn and C.~R.~Johnson.
  \emph{Matrix Analysis}, 2nd ed.\ Cambridge University Press, 2012.
\end{thebibliography}

\end{document}